% =========================================================
% ppt/src/minggu-07.tex
% Slide pertemuan 7
% Bab IX — Direktif Assembler dan Alur Translasi Program
%
% Sub-CPMK 2.3 | CPMK02 | CPL02
% =========================================================

\documentclass[aspectratio=169,11pt]{beamer}
% =========================================================
% ppt/src/preamble-ppt.tex
% Tema, warna, paket, dan makro bersama seluruh deck slide.
%
% Dipakai oleh setiap berkas minggu-NN.tex dengan pola:
%   \documentclass[aspectratio=169,11pt]{beamer}
%   % =========================================================
% ppt/src/preamble-ppt.tex
% Tema, warna, paket, dan makro bersama seluruh deck slide.
%
% Dipakai oleh setiap berkas minggu-NN.tex dengan pola:
%   \documentclass[aspectratio=169,11pt]{beamer}
%   % =========================================================
% ppt/src/preamble-ppt.tex
% Tema, warna, paket, dan makro bersama seluruh deck slide.
%
% Dipakai oleh setiap berkas minggu-NN.tex dengan pola:
%   \documentclass[aspectratio=169,11pt]{beamer}
%   \input{preamble-ppt}
%
% Palet warna disalin dari buku_ajar/config/styles.tex agar
% identitas visual slide sama dengan buku ajar.
% =========================================================

% ---------------------------------------------------------
% 1. Paket
% ---------------------------------------------------------
\usepackage[T1]{fontenc}
\usepackage[utf8]{inputenc}
\usepackage[indonesian]{babel}
\usepackage{lmodern}
\usepackage{tikz}
\usetikzlibrary{arrows.meta, positioning, calc, shapes.geometric}
\usepackage{booktabs}
\usepackage{array}
\usepackage{amsmath}
\usepackage{amsfonts}
\usepackage{pifont}          % \ding untuk tanda checklist
\usepackage{listings}

% ---------------------------------------------------------
% 2. Palet warna OBE (identik dengan buku)
% ---------------------------------------------------------
\definecolor{obeBlue}{HTML}{1F3864}
\definecolor{obeBlueLight}{HTML}{DCE4F5}
\definecolor{obeTeal}{HTML}{0F6E6E}
\definecolor{obeTealLight}{HTML}{D9F0EE}
\definecolor{obeGreen}{HTML}{2E6B34}
\definecolor{obeGreenLight}{HTML}{DEEFDF}
\definecolor{obeAmber}{HTML}{8A5A00}
\definecolor{obeAmberLight}{HTML}{FDF1D6}
\definecolor{obeMaroon}{HTML}{7A2233}
\definecolor{obeMaroonLight}{HTML}{F7E1E5}
\definecolor{obeGray}{HTML}{4A4A4A}
\definecolor{obeGrayLight}{HTML}{F0F0F0}
\definecolor{codeBack}{HTML}{F7F7F4}
\definecolor{codeFrame}{HTML}{C9C9C0}

% ---------------------------------------------------------
% 3. Identitas mata kuliah (dipakai seluruh deck)
% ---------------------------------------------------------
\newcommand{\NamaMK}{Bahasa Pemrograman Tingkat Rendah}
\newcommand{\NamaMKShort}{Bhs.\ Pemrograman Tingkat Rendah}
\newcommand{\KodeMK}{MWCA010}
\newcommand{\Prodi}{Program Studi Teknik Informatika}
\newcommand{\Fakultas}{Fakultas Teknologi Informasi}
\newcommand{\PerguruanTinggi}{Universitas Bale Bandung}
\newcommand{\DosenPengampu}{Cecep Suwanda, S.Si., M.Kom.}
\newcommand{\SemesterTeks}{Semester Ganjil}

% \PertemuanTeks dan \JudulBab ditetapkan di masing-masing
% berkas minggu-NN.tex sebelum \begin{document}.
\providecommand{\PertemuanTeks}{Pertemuan}

% ---------------------------------------------------------
% 4. Tema beamer
% ---------------------------------------------------------
\usetheme{default}
\setbeamertemplate{navigation symbols}{}

\setbeamercolor{normal text}{fg=black!88}
\setbeamercolor{structure}{fg=obeBlue}
\setbeamercolor{alerted text}{fg=obeMaroon}

% --- judul slide: pita biru penuh lebar ---
\setbeamercolor{frametitle}{fg=white,bg=obeBlue}
\setbeamerfont{frametitle}{series=\bfseries,size=\large}
\setbeamertemplate{frametitle}{%
    \nointerlineskip%
    \begin{beamercolorbox}[%
        wd=\paperwidth,ht=2.9ex,dp=1.5ex,
        leftskip=0.4cm,rightskip=0.4cm]{frametitle}%
        \usebeamerfont{frametitle}\insertframetitle%
    \end{beamercolorbox}%
    \nointerlineskip%
    \hbox{\textcolor{obeTeal}{\rule{\paperwidth}{1.1pt}}}%
}

% --- kaki slide: identitas mata kuliah, pertemuan, nomor ---
\setbeamercolor{footline}{fg=white,bg=obeGray}
\setbeamerfont{footline}{size=\scriptsize,family=\sffamily}
\setbeamertemplate{footline}{%
    \leavevmode%
    \hbox{%
        \begin{beamercolorbox}[wd=0.42\paperwidth,ht=2.1ex,dp=1.1ex,
            leftskip=0.4cm]{footline}%
            \usebeamerfont{footline}%
            \KodeMK\ \textbullet\ \NamaMKShort%
        \end{beamercolorbox}%
        \begin{beamercolorbox}[wd=0.32\paperwidth,ht=2.1ex,dp=1.1ex,
            center]{footline}%
            \usebeamerfont{footline}\PertemuanTeks%
        \end{beamercolorbox}%
        \begin{beamercolorbox}[wd=0.26\paperwidth,ht=2.1ex,dp=1.1ex,
            rightskip=0.4cm]{footline}%
            \usebeamerfont{footline}%
            \insertframenumber\,/\,\inserttotalframenumber%
        \end{beamercolorbox}%
    }%
    \vskip0pt%
}

% --- daftar ---
\setbeamercolor{itemize item}{fg=obeTeal}
\setbeamercolor{itemize subitem}{fg=obeTeal}
\setbeamercolor{enumerate item}{fg=obeBlue}
\setbeamertemplate{itemize item}{\raise0.2ex\hbox{\tiny$\blacksquare$}}
\setbeamertemplate{itemize subitem}{\raise0.2ex\hbox{\scriptsize$\bullet$}}

% --- kotak (pengganti tcolorbox pada buku) ---
\setbeamercolor{block title}{bg=obeBlueLight,fg=obeBlue}
\setbeamercolor{block body}{bg=obeBlueLight!45,fg=black!88}
\setbeamercolor{block title alerted}{bg=obeMaroonLight,fg=obeMaroon}
\setbeamercolor{block body alerted}{bg=obeMaroonLight!45,fg=black!88}
\setbeamercolor{block title example}{bg=obeTealLight,fg=obeTeal}
\setbeamercolor{block body example}{bg=obeTealLight!45,fg=black!88}

% ---------------------------------------------------------
% 5. Makro notasi (sama persis dengan buku)
%
% Catatan: \note TIDAK didefinisikan di sini karena beamer
% sudah memakai \note untuk catatan pembicara. Untuk catatan
% di slide dipakai \catatan.
% ---------------------------------------------------------
\newcommand{\code}[1]{\texttt{#1}}
\newcommand{\asm}[1]{\texttt{#1}}
\newcommand{\reg}[1]{\texttt{#1}}
\newcommand{\flag}[1]{\texttt{#1}}
\newcommand{\dir}[1]{\texttt{#1}}
\newcommand{\file}[1]{\texttt{#1}}
\newcommand{\eng}[1]{\textit{#1}}

\newcommand{\hex}[1]{\texttt{#1h}}
\newcommand{\bin}[1]{\texttt{#1b}}
\newcommand{\dec}[1]{\texttt{#1}}

\newcommand{\important}[1]{\textbf{Penting:} #1}
\newcommand{\catatan}[1]{\textbf{Catatan:} #1}
\newcommand{\term}[1]{\textbf{#1}}

\newcommand{\CPL}[1]{\textbf{CPL#1}}
\newcommand{\CPMK}[1]{\textbf{CPMK#1}}
\newcommand{\SubCPMK}[1]{\textbf{Sub-CPMK #1}}

\newcommand{\checkbox}{\ding{113}\hspace{0.35em}}
\newcommand{\checked}{\ding{51}\hspace{0.35em}}
\newcommand{\crossed}{\ding{55}\hspace{0.35em}}

\newcommand{\flowarrow}{\ensuremath{\rightarrow}}

% ---------------------------------------------------------
% 6. Makro tata letak slide
% ---------------------------------------------------------

% Keterangan sumber materi (ketertelusuran ke buku)
\newcommand{\Sumber}[1]{%
    \vfill
    {\scriptsize\color{obeGray}Sumber: #1\par}%
}

% Frame judul deck: \JudulDeck{<judul bab>}{<pertemuan>}
\newcommand{\JudulDeck}[2]{%
\begin{frame}[plain]
    \vspace{0.25cm}
    {\color{obeBlue}\sffamily\bfseries\small \NamaMK\ (\KodeMK)\par}
    \vspace{0.12cm}
    {\color{obeGray}\sffamily\scriptsize
     \Prodi\ \textbullet\ \Fakultas\ \textbullet\ \PerguruanTinggi\par}
    \vfill
    {\color{obeBlue}\sffamily\LARGE\bfseries #1\par}
    \vspace{0.4cm}
    {\color{obeTeal}\sffamily\large #2\par}
    \vspace{0.45cm}
    {\color{obeTeal}\rule{\textwidth}{1pt}}
    \vfill
    {\color{obeGray}\sffamily\small Dosen Pengampu\par}
    \vspace{0.08cm}
    {\sffamily\bfseries \DosenPengampu\par}
    \vspace{0.2cm}
    {\color{obeGray}\sffamily\small \SemesterTeks\par}
    \vspace{0.25cm}
\end{frame}
}

% Frame pemisah bagian
\newcommand{\PemisahBagian}[1]{%
\begin{frame}[plain]
    \vfill
    \begin{center}
        {\color{obeTeal}\sffamily\small\bfseries BAGIAN\par}
        \vspace{0.35cm}
        {\color{obeBlue}\sffamily\Large\bfseries #1\par}
        \vspace{0.35cm}
        {\color{obeTeal}\rule{0.55\textwidth}{0.9pt}}
    \end{center}
    \vfill
\end{frame}
}

% ---------------------------------------------------------
% 7. Gaya penampil kode
% ---------------------------------------------------------
\input{listing-ppt}

%
% Palet warna disalin dari buku_ajar/config/styles.tex agar
% identitas visual slide sama dengan buku ajar.
% =========================================================

% ---------------------------------------------------------
% 1. Paket
% ---------------------------------------------------------
\usepackage[T1]{fontenc}
\usepackage[utf8]{inputenc}
\usepackage[indonesian]{babel}
\usepackage{lmodern}
\usepackage{tikz}
\usetikzlibrary{arrows.meta, positioning, calc, shapes.geometric}
\usepackage{booktabs}
\usepackage{array}
\usepackage{amsmath}
\usepackage{amsfonts}
\usepackage{pifont}          % \ding untuk tanda checklist
\usepackage{listings}

% ---------------------------------------------------------
% 2. Palet warna OBE (identik dengan buku)
% ---------------------------------------------------------
\definecolor{obeBlue}{HTML}{1F3864}
\definecolor{obeBlueLight}{HTML}{DCE4F5}
\definecolor{obeTeal}{HTML}{0F6E6E}
\definecolor{obeTealLight}{HTML}{D9F0EE}
\definecolor{obeGreen}{HTML}{2E6B34}
\definecolor{obeGreenLight}{HTML}{DEEFDF}
\definecolor{obeAmber}{HTML}{8A5A00}
\definecolor{obeAmberLight}{HTML}{FDF1D6}
\definecolor{obeMaroon}{HTML}{7A2233}
\definecolor{obeMaroonLight}{HTML}{F7E1E5}
\definecolor{obeGray}{HTML}{4A4A4A}
\definecolor{obeGrayLight}{HTML}{F0F0F0}
\definecolor{codeBack}{HTML}{F7F7F4}
\definecolor{codeFrame}{HTML}{C9C9C0}

% ---------------------------------------------------------
% 3. Identitas mata kuliah (dipakai seluruh deck)
% ---------------------------------------------------------
\newcommand{\NamaMK}{Bahasa Pemrograman Tingkat Rendah}
\newcommand{\NamaMKShort}{Bhs.\ Pemrograman Tingkat Rendah}
\newcommand{\KodeMK}{MWCA010}
\newcommand{\Prodi}{Program Studi Teknik Informatika}
\newcommand{\Fakultas}{Fakultas Teknologi Informasi}
\newcommand{\PerguruanTinggi}{Universitas Bale Bandung}
\newcommand{\DosenPengampu}{Cecep Suwanda, S.Si., M.Kom.}
\newcommand{\SemesterTeks}{Semester Ganjil}

% \PertemuanTeks dan \JudulBab ditetapkan di masing-masing
% berkas minggu-NN.tex sebelum \begin{document}.
\providecommand{\PertemuanTeks}{Pertemuan}

% ---------------------------------------------------------
% 4. Tema beamer
% ---------------------------------------------------------
\usetheme{default}
\setbeamertemplate{navigation symbols}{}

\setbeamercolor{normal text}{fg=black!88}
\setbeamercolor{structure}{fg=obeBlue}
\setbeamercolor{alerted text}{fg=obeMaroon}

% --- judul slide: pita biru penuh lebar ---
\setbeamercolor{frametitle}{fg=white,bg=obeBlue}
\setbeamerfont{frametitle}{series=\bfseries,size=\large}
\setbeamertemplate{frametitle}{%
    \nointerlineskip%
    \begin{beamercolorbox}[%
        wd=\paperwidth,ht=2.9ex,dp=1.5ex,
        leftskip=0.4cm,rightskip=0.4cm]{frametitle}%
        \usebeamerfont{frametitle}\insertframetitle%
    \end{beamercolorbox}%
    \nointerlineskip%
    \hbox{\textcolor{obeTeal}{\rule{\paperwidth}{1.1pt}}}%
}

% --- kaki slide: identitas mata kuliah, pertemuan, nomor ---
\setbeamercolor{footline}{fg=white,bg=obeGray}
\setbeamerfont{footline}{size=\scriptsize,family=\sffamily}
\setbeamertemplate{footline}{%
    \leavevmode%
    \hbox{%
        \begin{beamercolorbox}[wd=0.42\paperwidth,ht=2.1ex,dp=1.1ex,
            leftskip=0.4cm]{footline}%
            \usebeamerfont{footline}%
            \KodeMK\ \textbullet\ \NamaMKShort%
        \end{beamercolorbox}%
        \begin{beamercolorbox}[wd=0.32\paperwidth,ht=2.1ex,dp=1.1ex,
            center]{footline}%
            \usebeamerfont{footline}\PertemuanTeks%
        \end{beamercolorbox}%
        \begin{beamercolorbox}[wd=0.26\paperwidth,ht=2.1ex,dp=1.1ex,
            rightskip=0.4cm]{footline}%
            \usebeamerfont{footline}%
            \insertframenumber\,/\,\inserttotalframenumber%
        \end{beamercolorbox}%
    }%
    \vskip0pt%
}

% --- daftar ---
\setbeamercolor{itemize item}{fg=obeTeal}
\setbeamercolor{itemize subitem}{fg=obeTeal}
\setbeamercolor{enumerate item}{fg=obeBlue}
\setbeamertemplate{itemize item}{\raise0.2ex\hbox{\tiny$\blacksquare$}}
\setbeamertemplate{itemize subitem}{\raise0.2ex\hbox{\scriptsize$\bullet$}}

% --- kotak (pengganti tcolorbox pada buku) ---
\setbeamercolor{block title}{bg=obeBlueLight,fg=obeBlue}
\setbeamercolor{block body}{bg=obeBlueLight!45,fg=black!88}
\setbeamercolor{block title alerted}{bg=obeMaroonLight,fg=obeMaroon}
\setbeamercolor{block body alerted}{bg=obeMaroonLight!45,fg=black!88}
\setbeamercolor{block title example}{bg=obeTealLight,fg=obeTeal}
\setbeamercolor{block body example}{bg=obeTealLight!45,fg=black!88}

% ---------------------------------------------------------
% 5. Makro notasi (sama persis dengan buku)
%
% Catatan: \note TIDAK didefinisikan di sini karena beamer
% sudah memakai \note untuk catatan pembicara. Untuk catatan
% di slide dipakai \catatan.
% ---------------------------------------------------------
\newcommand{\code}[1]{\texttt{#1}}
\newcommand{\asm}[1]{\texttt{#1}}
\newcommand{\reg}[1]{\texttt{#1}}
\newcommand{\flag}[1]{\texttt{#1}}
\newcommand{\dir}[1]{\texttt{#1}}
\newcommand{\file}[1]{\texttt{#1}}
\newcommand{\eng}[1]{\textit{#1}}

\newcommand{\hex}[1]{\texttt{#1h}}
\newcommand{\bin}[1]{\texttt{#1b}}
\newcommand{\dec}[1]{\texttt{#1}}

\newcommand{\important}[1]{\textbf{Penting:} #1}
\newcommand{\catatan}[1]{\textbf{Catatan:} #1}
\newcommand{\term}[1]{\textbf{#1}}

\newcommand{\CPL}[1]{\textbf{CPL#1}}
\newcommand{\CPMK}[1]{\textbf{CPMK#1}}
\newcommand{\SubCPMK}[1]{\textbf{Sub-CPMK #1}}

\newcommand{\checkbox}{\ding{113}\hspace{0.35em}}
\newcommand{\checked}{\ding{51}\hspace{0.35em}}
\newcommand{\crossed}{\ding{55}\hspace{0.35em}}

\newcommand{\flowarrow}{\ensuremath{\rightarrow}}

% ---------------------------------------------------------
% 6. Makro tata letak slide
% ---------------------------------------------------------

% Keterangan sumber materi (ketertelusuran ke buku)
\newcommand{\Sumber}[1]{%
    \vfill
    {\scriptsize\color{obeGray}Sumber: #1\par}%
}

% Frame judul deck: \JudulDeck{<judul bab>}{<pertemuan>}
\newcommand{\JudulDeck}[2]{%
\begin{frame}[plain]
    \vspace{0.25cm}
    {\color{obeBlue}\sffamily\bfseries\small \NamaMK\ (\KodeMK)\par}
    \vspace{0.12cm}
    {\color{obeGray}\sffamily\scriptsize
     \Prodi\ \textbullet\ \Fakultas\ \textbullet\ \PerguruanTinggi\par}
    \vfill
    {\color{obeBlue}\sffamily\LARGE\bfseries #1\par}
    \vspace{0.4cm}
    {\color{obeTeal}\sffamily\large #2\par}
    \vspace{0.45cm}
    {\color{obeTeal}\rule{\textwidth}{1pt}}
    \vfill
    {\color{obeGray}\sffamily\small Dosen Pengampu\par}
    \vspace{0.08cm}
    {\sffamily\bfseries \DosenPengampu\par}
    \vspace{0.2cm}
    {\color{obeGray}\sffamily\small \SemesterTeks\par}
    \vspace{0.25cm}
\end{frame}
}

% Frame pemisah bagian
\newcommand{\PemisahBagian}[1]{%
\begin{frame}[plain]
    \vfill
    \begin{center}
        {\color{obeTeal}\sffamily\small\bfseries BAGIAN\par}
        \vspace{0.35cm}
        {\color{obeBlue}\sffamily\Large\bfseries #1\par}
        \vspace{0.35cm}
        {\color{obeTeal}\rule{0.55\textwidth}{0.9pt}}
    \end{center}
    \vfill
\end{frame}
}

% ---------------------------------------------------------
% 7. Gaya penampil kode
% ---------------------------------------------------------
% =========================================================
% ppt/src/listing-ppt.tex
% Gaya penampil kode untuk slide.
%
% Diturunkan dari buku_ajar/config/listings.tex, tetapi
% disesuaikan untuk slide: ukuran lebih kecil (karena frame
% 16:9 dan kode harus utuh terlihat) dan jarak lebih rapat.
% =========================================================

% ---------------------------------------------------------
% 1. Definisi bahasa MASM/TASM 8086 (salinan dari buku)
% ---------------------------------------------------------
\lstdefinelanguage{MASM}{
    morekeywords={
        MOV, XCHG, LEA, LDS, LES, PUSH, POP, PUSHA, POPA, PUSHF, POPF,
        LAHF, SAHF, XLAT, IN, OUT, CBW, CWD,
        ADD, ADC, SUB, SBB, INC, DEC, NEG, CMP, MUL, IMUL, DIV, IDIV,
        AAA, AAS, AAM, AAD, DAA, DAS,
        AND, OR, XOR, NOT, TEST, SHL, SHR, SAL, SAR, ROL, ROR, RCL, RCR,
        JMP, JE, JZ, JNE, JNZ, JG, JGE, JL, JLE, JA, JAE, JB, JBE,
        JC, JNC, JS, JNS, JO, JNO, JP, JPE, JNP, JPO, JCXZ,
        LOOP, LOOPE, LOOPZ, LOOPNE, LOOPNZ, CALL, RET, INT, IRET, INTO,
        HLT, NOP, CLI, STI, CLD, STD, CLC, STC, CMC, WAIT, ESC, LOCK,
        DB, DW, DD, DQ, DT, EQU, ORG, SEGMENT, ENDS, ASSUME, END,
        PROC, ENDP, MACRO, ENDM, LOCAL, PUBLIC, EXTRN, EXTERN,
        INCLUDE, INCLUDELIB, MODEL, STACK, CODE, DATA, CONST, ALIGN, EVEN,
        OFFSET, PTR, BYTE, WORD, DWORD, QWORD, TBYTE,
        NEAR, FAR, SHORT, LABEL, TYPE, SIZE, LENGTH, DUP, THIS,
        AX, BX, CX, DX, AL, AH, BL, BH, CL, CH, DL, DH,
        SP, BP, SI, DI, IP, CS, DS, SS, ES, FLAGS,
        CF, PF, AF, ZF, TF, IF, DF, OF
    },
    sensitive=false,
    morecomment=[l]{;},
    morestring=[b]',
    morestring=[b]"
}

% ---------------------------------------------------------
% 2. Gaya kode assembly untuk slide
% ---------------------------------------------------------
\lstdefinestyle{asmslide}{
    language=MASM,
    basicstyle=\ttfamily\scriptsize,
    keywordstyle=\color{obeBlue}\bfseries,
    commentstyle=\color{obeGreen}\itshape,
    stringstyle=\color{obeMaroon},
    numberstyle=\tiny\color{obeGray},
    numbers=left,
    numbersep=5pt,
    stepnumber=1,
    backgroundcolor=\color{codeBack},
    frame=single,
    rulecolor=\color{codeFrame},
    framesep=4pt,
    breaklines=true,
    breakatwhitespace=false,
    % CATATAN PENTING: opsi `postbreak` TIDAK boleh dipakai di sini.
    % Di beamer, `breaklines=true` bersama `postbreak=<kode>` membuat
    % listings gagal dirakit dengan galat
    %   "Missing number, treated as zero <to be read again> \kern".
    % Karena \lstset{style=asmslide} bersifat global, opsi postbreak
    % juga bocor ke gaya lain (mis. outslide/keluaran) dan merusaknya.
    % Di buku ajar (kelas `book`) justru sebaliknya: postbreak aman dan
    % dipakai. Jadi larangan ini khusus untuk slide.
    showstringspaces=false,
    tabsize=4,
    columns=fullflexible,
    keepspaces=true,
    upquote=true,
    xleftmargin=12pt,
    aboveskip=0.4em,
    belowskip=0.4em
}

% ---------------------------------------------------------
% 3. Gaya keluaran program / sesi debugger untuk slide
% ---------------------------------------------------------
\lstdefinestyle{outslide}{
    basicstyle=\ttfamily\scriptsize,
    keywordstyle=\color{black},
    commentstyle=\color{obeGray},
    stringstyle=\color{black},
    numbers=none,
    backgroundcolor=\color{obeGrayLight},
    frame=single,
    rulecolor=\color{codeFrame},
    framesep=4pt,
    breaklines=true,
    showstringspaces=false,
    columns=fullflexible,
    keepspaces=true,
    upquote=true,
    xleftmargin=4pt,
    aboveskip=0.4em,
    belowskip=0.4em
}

\lstset{style=asmslide}

% ---------------------------------------------------------
% 4. Environment ringkas untuk dipakai di dalam frame
%
% Frame yang memuat kedua environment ini WAJIB memakai
% opsi [fragile], karena isinya verbatim.
% ---------------------------------------------------------
\lstnewenvironment{kodeasm}[1][]{%
    \lstset{style=asmslide, #1}%
}{}

\lstnewenvironment{keluaran}[1][]{%
    \lstset{style=outslide, #1}%
}{}


%
% Palet warna disalin dari buku_ajar/config/styles.tex agar
% identitas visual slide sama dengan buku ajar.
% =========================================================

% ---------------------------------------------------------
% 1. Paket
% ---------------------------------------------------------
\usepackage[T1]{fontenc}
\usepackage[utf8]{inputenc}
\usepackage[indonesian]{babel}
\usepackage{lmodern}
\usepackage{tikz}
\usetikzlibrary{arrows.meta, positioning, calc, shapes.geometric}
\usepackage{booktabs}
\usepackage{array}
\usepackage{amsmath}
\usepackage{amsfonts}
\usepackage{pifont}          % \ding untuk tanda checklist
\usepackage{listings}

% ---------------------------------------------------------
% 2. Palet warna OBE (identik dengan buku)
% ---------------------------------------------------------
\definecolor{obeBlue}{HTML}{1F3864}
\definecolor{obeBlueLight}{HTML}{DCE4F5}
\definecolor{obeTeal}{HTML}{0F6E6E}
\definecolor{obeTealLight}{HTML}{D9F0EE}
\definecolor{obeGreen}{HTML}{2E6B34}
\definecolor{obeGreenLight}{HTML}{DEEFDF}
\definecolor{obeAmber}{HTML}{8A5A00}
\definecolor{obeAmberLight}{HTML}{FDF1D6}
\definecolor{obeMaroon}{HTML}{7A2233}
\definecolor{obeMaroonLight}{HTML}{F7E1E5}
\definecolor{obeGray}{HTML}{4A4A4A}
\definecolor{obeGrayLight}{HTML}{F0F0F0}
\definecolor{codeBack}{HTML}{F7F7F4}
\definecolor{codeFrame}{HTML}{C9C9C0}

% ---------------------------------------------------------
% 3. Identitas mata kuliah (dipakai seluruh deck)
% ---------------------------------------------------------
\newcommand{\NamaMK}{Bahasa Pemrograman Tingkat Rendah}
\newcommand{\NamaMKShort}{Bhs.\ Pemrograman Tingkat Rendah}
\newcommand{\KodeMK}{MWCA010}
\newcommand{\Prodi}{Program Studi Teknik Informatika}
\newcommand{\Fakultas}{Fakultas Teknologi Informasi}
\newcommand{\PerguruanTinggi}{Universitas Bale Bandung}
\newcommand{\DosenPengampu}{Cecep Suwanda, S.Si., M.Kom.}
\newcommand{\SemesterTeks}{Semester Ganjil}

% \PertemuanTeks dan \JudulBab ditetapkan di masing-masing
% berkas minggu-NN.tex sebelum \begin{document}.
\providecommand{\PertemuanTeks}{Pertemuan}

% ---------------------------------------------------------
% 4. Tema beamer
% ---------------------------------------------------------
\usetheme{default}
\setbeamertemplate{navigation symbols}{}

\setbeamercolor{normal text}{fg=black!88}
\setbeamercolor{structure}{fg=obeBlue}
\setbeamercolor{alerted text}{fg=obeMaroon}

% --- judul slide: pita biru penuh lebar ---
\setbeamercolor{frametitle}{fg=white,bg=obeBlue}
\setbeamerfont{frametitle}{series=\bfseries,size=\large}
\setbeamertemplate{frametitle}{%
    \nointerlineskip%
    \begin{beamercolorbox}[%
        wd=\paperwidth,ht=2.9ex,dp=1.5ex,
        leftskip=0.4cm,rightskip=0.4cm]{frametitle}%
        \usebeamerfont{frametitle}\insertframetitle%
    \end{beamercolorbox}%
    \nointerlineskip%
    \hbox{\textcolor{obeTeal}{\rule{\paperwidth}{1.1pt}}}%
}

% --- kaki slide: identitas mata kuliah, pertemuan, nomor ---
\setbeamercolor{footline}{fg=white,bg=obeGray}
\setbeamerfont{footline}{size=\scriptsize,family=\sffamily}
\setbeamertemplate{footline}{%
    \leavevmode%
    \hbox{%
        \begin{beamercolorbox}[wd=0.42\paperwidth,ht=2.1ex,dp=1.1ex,
            leftskip=0.4cm]{footline}%
            \usebeamerfont{footline}%
            \KodeMK\ \textbullet\ \NamaMKShort%
        \end{beamercolorbox}%
        \begin{beamercolorbox}[wd=0.32\paperwidth,ht=2.1ex,dp=1.1ex,
            center]{footline}%
            \usebeamerfont{footline}\PertemuanTeks%
        \end{beamercolorbox}%
        \begin{beamercolorbox}[wd=0.26\paperwidth,ht=2.1ex,dp=1.1ex,
            rightskip=0.4cm]{footline}%
            \usebeamerfont{footline}%
            \insertframenumber\,/\,\inserttotalframenumber%
        \end{beamercolorbox}%
    }%
    \vskip0pt%
}

% --- daftar ---
\setbeamercolor{itemize item}{fg=obeTeal}
\setbeamercolor{itemize subitem}{fg=obeTeal}
\setbeamercolor{enumerate item}{fg=obeBlue}
\setbeamertemplate{itemize item}{\raise0.2ex\hbox{\tiny$\blacksquare$}}
\setbeamertemplate{itemize subitem}{\raise0.2ex\hbox{\scriptsize$\bullet$}}

% --- kotak (pengganti tcolorbox pada buku) ---
\setbeamercolor{block title}{bg=obeBlueLight,fg=obeBlue}
\setbeamercolor{block body}{bg=obeBlueLight!45,fg=black!88}
\setbeamercolor{block title alerted}{bg=obeMaroonLight,fg=obeMaroon}
\setbeamercolor{block body alerted}{bg=obeMaroonLight!45,fg=black!88}
\setbeamercolor{block title example}{bg=obeTealLight,fg=obeTeal}
\setbeamercolor{block body example}{bg=obeTealLight!45,fg=black!88}

% ---------------------------------------------------------
% 5. Makro notasi (sama persis dengan buku)
%
% Catatan: \note TIDAK didefinisikan di sini karena beamer
% sudah memakai \note untuk catatan pembicara. Untuk catatan
% di slide dipakai \catatan.
% ---------------------------------------------------------
\newcommand{\code}[1]{\texttt{#1}}
\newcommand{\asm}[1]{\texttt{#1}}
\newcommand{\reg}[1]{\texttt{#1}}
\newcommand{\flag}[1]{\texttt{#1}}
\newcommand{\dir}[1]{\texttt{#1}}
\newcommand{\file}[1]{\texttt{#1}}
\newcommand{\eng}[1]{\textit{#1}}

\newcommand{\hex}[1]{\texttt{#1h}}
\newcommand{\bin}[1]{\texttt{#1b}}
\newcommand{\dec}[1]{\texttt{#1}}

\newcommand{\important}[1]{\textbf{Penting:} #1}
\newcommand{\catatan}[1]{\textbf{Catatan:} #1}
\newcommand{\term}[1]{\textbf{#1}}

\newcommand{\CPL}[1]{\textbf{CPL#1}}
\newcommand{\CPMK}[1]{\textbf{CPMK#1}}
\newcommand{\SubCPMK}[1]{\textbf{Sub-CPMK #1}}

\newcommand{\checkbox}{\ding{113}\hspace{0.35em}}
\newcommand{\checked}{\ding{51}\hspace{0.35em}}
\newcommand{\crossed}{\ding{55}\hspace{0.35em}}

\newcommand{\flowarrow}{\ensuremath{\rightarrow}}

% ---------------------------------------------------------
% 6. Makro tata letak slide
% ---------------------------------------------------------

% Keterangan sumber materi (ketertelusuran ke buku)
\newcommand{\Sumber}[1]{%
    \vfill
    {\scriptsize\color{obeGray}Sumber: #1\par}%
}

% Frame judul deck: \JudulDeck{<judul bab>}{<pertemuan>}
\newcommand{\JudulDeck}[2]{%
\begin{frame}[plain]
    \vspace{0.25cm}
    {\color{obeBlue}\sffamily\bfseries\small \NamaMK\ (\KodeMK)\par}
    \vspace{0.12cm}
    {\color{obeGray}\sffamily\scriptsize
     \Prodi\ \textbullet\ \Fakultas\ \textbullet\ \PerguruanTinggi\par}
    \vfill
    {\color{obeBlue}\sffamily\LARGE\bfseries #1\par}
    \vspace{0.4cm}
    {\color{obeTeal}\sffamily\large #2\par}
    \vspace{0.45cm}
    {\color{obeTeal}\rule{\textwidth}{1pt}}
    \vfill
    {\color{obeGray}\sffamily\small Dosen Pengampu\par}
    \vspace{0.08cm}
    {\sffamily\bfseries \DosenPengampu\par}
    \vspace{0.2cm}
    {\color{obeGray}\sffamily\small \SemesterTeks\par}
    \vspace{0.25cm}
\end{frame}
}

% Frame pemisah bagian
\newcommand{\PemisahBagian}[1]{%
\begin{frame}[plain]
    \vfill
    \begin{center}
        {\color{obeTeal}\sffamily\small\bfseries BAGIAN\par}
        \vspace{0.35cm}
        {\color{obeBlue}\sffamily\Large\bfseries #1\par}
        \vspace{0.35cm}
        {\color{obeTeal}\rule{0.55\textwidth}{0.9pt}}
    \end{center}
    \vfill
\end{frame}
}

% ---------------------------------------------------------
% 7. Gaya penampil kode
% ---------------------------------------------------------
% =========================================================
% ppt/src/listing-ppt.tex
% Gaya penampil kode untuk slide.
%
% Diturunkan dari buku_ajar/config/listings.tex, tetapi
% disesuaikan untuk slide: ukuran lebih kecil (karena frame
% 16:9 dan kode harus utuh terlihat) dan jarak lebih rapat.
% =========================================================

% ---------------------------------------------------------
% 1. Definisi bahasa MASM/TASM 8086 (salinan dari buku)
% ---------------------------------------------------------
\lstdefinelanguage{MASM}{
    morekeywords={
        MOV, XCHG, LEA, LDS, LES, PUSH, POP, PUSHA, POPA, PUSHF, POPF,
        LAHF, SAHF, XLAT, IN, OUT, CBW, CWD,
        ADD, ADC, SUB, SBB, INC, DEC, NEG, CMP, MUL, IMUL, DIV, IDIV,
        AAA, AAS, AAM, AAD, DAA, DAS,
        AND, OR, XOR, NOT, TEST, SHL, SHR, SAL, SAR, ROL, ROR, RCL, RCR,
        JMP, JE, JZ, JNE, JNZ, JG, JGE, JL, JLE, JA, JAE, JB, JBE,
        JC, JNC, JS, JNS, JO, JNO, JP, JPE, JNP, JPO, JCXZ,
        LOOP, LOOPE, LOOPZ, LOOPNE, LOOPNZ, CALL, RET, INT, IRET, INTO,
        HLT, NOP, CLI, STI, CLD, STD, CLC, STC, CMC, WAIT, ESC, LOCK,
        DB, DW, DD, DQ, DT, EQU, ORG, SEGMENT, ENDS, ASSUME, END,
        PROC, ENDP, MACRO, ENDM, LOCAL, PUBLIC, EXTRN, EXTERN,
        INCLUDE, INCLUDELIB, MODEL, STACK, CODE, DATA, CONST, ALIGN, EVEN,
        OFFSET, PTR, BYTE, WORD, DWORD, QWORD, TBYTE,
        NEAR, FAR, SHORT, LABEL, TYPE, SIZE, LENGTH, DUP, THIS,
        AX, BX, CX, DX, AL, AH, BL, BH, CL, CH, DL, DH,
        SP, BP, SI, DI, IP, CS, DS, SS, ES, FLAGS,
        CF, PF, AF, ZF, TF, IF, DF, OF
    },
    sensitive=false,
    morecomment=[l]{;},
    morestring=[b]',
    morestring=[b]"
}

% ---------------------------------------------------------
% 2. Gaya kode assembly untuk slide
% ---------------------------------------------------------
\lstdefinestyle{asmslide}{
    language=MASM,
    basicstyle=\ttfamily\scriptsize,
    keywordstyle=\color{obeBlue}\bfseries,
    commentstyle=\color{obeGreen}\itshape,
    stringstyle=\color{obeMaroon},
    numberstyle=\tiny\color{obeGray},
    numbers=left,
    numbersep=5pt,
    stepnumber=1,
    backgroundcolor=\color{codeBack},
    frame=single,
    rulecolor=\color{codeFrame},
    framesep=4pt,
    breaklines=true,
    breakatwhitespace=false,
    % CATATAN PENTING: opsi `postbreak` TIDAK boleh dipakai di sini.
    % Di beamer, `breaklines=true` bersama `postbreak=<kode>` membuat
    % listings gagal dirakit dengan galat
    %   "Missing number, treated as zero <to be read again> \kern".
    % Karena \lstset{style=asmslide} bersifat global, opsi postbreak
    % juga bocor ke gaya lain (mis. outslide/keluaran) dan merusaknya.
    % Di buku ajar (kelas `book`) justru sebaliknya: postbreak aman dan
    % dipakai. Jadi larangan ini khusus untuk slide.
    showstringspaces=false,
    tabsize=4,
    columns=fullflexible,
    keepspaces=true,
    upquote=true,
    xleftmargin=12pt,
    aboveskip=0.4em,
    belowskip=0.4em
}

% ---------------------------------------------------------
% 3. Gaya keluaran program / sesi debugger untuk slide
% ---------------------------------------------------------
\lstdefinestyle{outslide}{
    basicstyle=\ttfamily\scriptsize,
    keywordstyle=\color{black},
    commentstyle=\color{obeGray},
    stringstyle=\color{black},
    numbers=none,
    backgroundcolor=\color{obeGrayLight},
    frame=single,
    rulecolor=\color{codeFrame},
    framesep=4pt,
    breaklines=true,
    showstringspaces=false,
    columns=fullflexible,
    keepspaces=true,
    upquote=true,
    xleftmargin=4pt,
    aboveskip=0.4em,
    belowskip=0.4em
}

\lstset{style=asmslide}

% ---------------------------------------------------------
% 4. Environment ringkas untuk dipakai di dalam frame
%
% Frame yang memuat kedua environment ini WAJIB memakai
% opsi [fragile], karena isinya verbatim.
% ---------------------------------------------------------
\lstnewenvironment{kodeasm}[1][]{%
    \lstset{style=asmslide, #1}%
}{}

\lstnewenvironment{keluaran}[1][]{%
    \lstset{style=outslide, #1}%
}{}



\renewcommand{\PertemuanTeks}{Pertemuan 7}

\begin{document}

% =========================================================
% 1. JUDUL
% =========================================================
\JudulDeck{Direktif Assembler\\ dan Alur Translasi Program}{Pertemuan 7}

% =========================================================
% 2. CAPAIAN DAN TUJUAN
% =========================================================
\begin{frame}{Capaian dan Tujuan Pembelajaran}
\begin{block}{Sub-CPMK 2.3}
\footnotesize Menggunakan direktif assembler (\dir{DB}, \dir{DW},
\dir{DD}, \dir{EQU}, \dir{ORG}, \dir{SEGMENT}) dan menjelaskan alur
translasi \eng{source} $\rightarrow$ \eng{object} $\rightarrow$
\eng{executable} oleh assembler dan linker.
\end{block}

\vspace{0.05em}
{\small\textbf{CPMK yang didukung:} \CPMK{02}\quad\textbullet\quad
\textbf{CPL yang didukung:} \CPL{02}\par}

\vspace{0.05em}
\begin{exampleblock}{Setelah pertemuan ini, mahasiswa mampu:}
\fontsize{9}{10}\selectfont
\begin{enumerate}
    \item membedakan instruksi mesin dan direktif assembler serta
          menjelaskan mengapa direktif tidak menghasilkan kode mesin;
    \item menggunakan direktif data \dir{DB}, \dir{DW}, \dir{DD},
          dan \dir{DT} beserta literal string, operator \dir{DUP},
          penanda \code{?}, dan pengaruh urutan \eng{little-endian};
    \item membedakan direktif \dir{EQU} dan \code{=} beserta
          akibatnya bila simbol didefinisikan ulang;
    \item menggunakan direktif lokasi \dir{ORG},
          \dir{SEGMENT}/\dir{ENDS}, \dir{ASSUME}, dan \dir{END}
          sesuai bentuk program \file{.COM} maupun \file{.EXE},
          serta memilih model memori \dir{.MODEL} yang tepat;
    \item menjelaskan tugas assembler (tabel simbol, \eng{offset})
          dan linker (\eng{relocation}, resolusi simbol), lalu
          menelusuri alur \file{.asm} $\rightarrow$ \file{.obj}
          $\rightarrow$ \file{.exe} dengan TASM dan TLINK.
\end{enumerate}
\end{exampleblock}
\end{frame}

% =========================================================
% 3. PETA MATERI
% =========================================================
\begin{frame}{Peta Materi}
\small
\begin{enumerate}
    \item \textbf{Instruksi Mesin, Direktif Assembler, dan Direktif
          Data}\\
          {\footnotesize dua jenis baris kode sumber; \dir{DB},
          \dir{DW}, \dir{DD}, \dir{DT}; \dir{DUP}; penanda \code{?};
          urutan \eng{little-endian}}

    \vspace{0.25em}
    \item \textbf{Direktif Simbol, Lokasi, dan Struktur Program}\\
          {\footnotesize \dir{EQU} dan \code{=}; \dir{ORG},
          \dir{SEGMENT}/\dir{ENDS}, \dir{ASSUME}, \dir{END};
          \dir{.MODEL}; \dir{PROC}/\dir{ENDP}}

    \vspace{0.25em}
    \item \textbf{Peran Assembler dan Linker dalam Alur Translasi}\\
          {\footnotesize dua lintasan assembler; tabel simbol;
          berkas \file{.lst}; resolusi simbol dan \eng{relocation};
          berkas \file{.map}; kesalahan perakitan}
\end{enumerate}

\vspace{0.4em}
\begin{block}{Praktik pada pertemuan ini}
\small Merakit \file{direktif-data.asm} dengan TASM dan TLINK, lalu
memeriksa berkas \file{.lst} dan \file{.map}.
\end{block}
\end{frame}

% =========================================================
% BAGIAN 1
% =========================================================
\PemisahBagian{Instruksi Mesin, Direktif Assembler, dan Direktif Data}

% ---------------------------------------------------------
\begin{frame}[fragile]{Dua Jenis Baris dalam Kode Sumber}
\footnotesize
\textbf{Instruksi mesin} --- mnemonic seperti \asm{MOV}, \asm{ADD},
dan \asm{INT} diterjemahkan menjadi \term{opcode}, yaitu kode mesin
yang benar-benar dieksekusi prosesor. Instruksi hanya bekerja pada
saat program berjalan (\eng{run time}).

\vspace{0.35em}
\textbf{Direktif assembler} --- hanya berlaku pada saat perakitan
(\eng{assemble time}); ia memberi tahu assembler bagaimana menata
program. Disebut \eng{pseudo-op} karena bentuknya mirip instruksi
tetapi \important{tidak menghasilkan kode mesin apa pun}.

\vspace{0.3em}
{\color{obeGray}Pada berkas \file{.lst}, baris instruksi memiliki
kolom kode mesin, sedangkan baris direktif kosong:\par}
\vspace{0.15em}
\begin{keluaran}
  4                             .DATA
  5  0000  48 61 6C 6F        pesan   DB 'Halo'
  6                             panjang EQU 5
  8                             .CODE
  9  0000  B8 ----            START:  MOV AX, @DATA
 10  0003  8E D8                      MOV DS, AX
\end{keluaran}
\Sumber{Bab IX, Dua Jenis Baris dalam Kode Sumber}
\end{frame}

% ---------------------------------------------------------
\begin{frame}{Direktif Data dan Ukuran yang Disediakannya}
\footnotesize
\begin{center}
\begin{tabular}{@{}p{1.5cm}p{3.2cm}p{3.6cm}p{4.2cm}@{}}
    \toprule
    \textbf{Direktif} & \textbf{Fungsi} & \textbf{Contoh} &
    \textbf{Ukuran / efek} \\
    \midrule
    \dir{DB} & Data 1 byte & \code{DB 0Ah} & 1 byte per nilai \\
    \addlinespace
    \dir{DW} & Data 2 byte & \code{DW 1234h} &
    2 byte; \eng{little-endian} \\
    \addlinespace
    \dir{DD} & Data 4 byte & \code{DD 12345678h} &
    4 byte; \eng{little-endian} \\
    \addlinespace
    \dir{DT} & Data 10 byte & \code{DT 1.0} &
    10 byte (bilangan pecahan) \\
    \addlinespace
    \dir{DUP} & Mengulang satu pola & \code{DB 10 DUP(0)} &
    10 byte bernilai nol \\
    \addlinespace
    \code{?} & Ruang tanpa nilai awal & \code{DB 8 DUP(?)} &
    8 byte tanpa isi \\
    \addlinespace
    \code{\$} & Penghitung lokasi & \code{EQU \$ - pesan} &
    Nilai \eng{offset} saat itu \\
    \bottomrule
\end{tabular}
\end{center}

\vspace{0.25em}
{\color{obeGray}Direktif yang mengisi memori (\dir{DB}, \dir{DW},
\dir{DD}) tetap mempengaruhi isi berkas \file{.exe}. Yang tidak
dihasilkannya adalah \emph{instruksi}.\par}
\Sumber{Bab IX, Tabel ringkasan direktif assembler}
\end{frame}

% ---------------------------------------------------------
\begin{frame}{Bentuk Penulisan yang Sering Dipakai}
\footnotesize
\begin{enumerate}
    \item \textbf{Literal string.} \code{DB 'Halo'} menyediakan empat
          byte: \hex{48}, \hex{61}, \hex{6C}, \hex{6F}. String tidak
          otomatis diakhiri apa pun, sehingga penanda akhir harus
          ditambahkan sendiri: \code{DB 'Halo', '\$'}.

    \vspace{0.3em}
    \item \textbf{Beberapa nilai sekaligus.} \code{DW 1, 2, 3}
          menyediakan tiga kata --- enam byte --- dengan urutan
          sesuai penulisan.

    \vspace{0.3em}
    \item \textbf{Operator \dir{DUP}.} Bentuk
          \code{<jumlah> DUP(<nilai>)} mengulang satu pola; bentuk
          bersarang juga diizinkan, misalnya
          \code{DB 2 DUP(3 DUP(0))}.

    \vspace{0.3em}
    \item \textbf{Penanda \code{?}.} Tanda tanya berarti sediakan
          ruang \emph{tanpa} nilai awal, seperti pada
          \code{buffer DB 8 DUP(?)}.

    \vspace{0.3em}
    \item \textbf{Penghitung lokasi \code{\$}.} Simbol \code{\$}
          berisi \eng{offset} baris yang sedang dirakit, sehingga
          \code{panjang EQU \$ - pesan} menghitung panjang string
          secara otomatis.
\end{enumerate}
\Sumber{Bab IX, Direktif Data}
\end{frame}

% ---------------------------------------------------------
\begin{frame}{Urutan \eng{Little-Endian} pada \dir{DW} dan \dir{DD}}
\small Prosesor 8086 menyimpan bilangan multi-byte dengan urutan
\term{little-endian}: \important{byte berpangkat rendah disimpan pada
alamat yang lebih kecil}.

\vspace{0.4em}
\begin{center}
\begin{tikzpicture}[
    sel/.style={draw=obeBlue, line width=0.7pt, minimum width=0.95cm,
        minimum height=0.62cm, align=center, font=\ttfamily\small,
        fill=obeBlueLight, inner sep=0pt},
    selD/.style={sel, draw=obeTeal, fill=obeTealLight},
    addr/.style={font=\scriptsize, text=obeGray},
    nama/.style={font=\small\bfseries, text=obeBlue, anchor=east},
    catat/.style={font=\scriptsize\itshape, text=obeMaroon}
]
    \node[nama] at (-0.55,0) {\dir{DW} \hex{1234}};
    \node[sel] (w0) at (0,0) {34};
    \node[sel, right=0pt of w0] (w1) {12};
    \node[addr, below=0.06cm of w0] {\hex{0300}};
    \node[addr, below=0.06cm of w1] {\hex{0301}};
    \node[catat, above=0.12cm of w1, xshift=-0.55cm]
        {byte rendah lebih dulu};

    \node[nama] at (-0.55,-2.0) {\dir{DD} \hex{12345678}};
    \node[selD] (d0) at (0,-2.0) {78};
    \node[selD, right=0pt of d0] (d1) {56};
    \node[selD, right=0pt of d1] (d2) {34};
    \node[selD, right=0pt of d2] (d3) {12};
    \node[addr, below=0.06cm of d0] {\hex{0302}};
    \node[addr, below=0.06cm of d1] {\hex{0303}};
    \node[addr, below=0.06cm of d2] {\hex{0304}};
    \node[addr, below=0.06cm of d3] {\hex{0305}};
\end{tikzpicture}
\end{center}

\vspace{0.3em}
{\footnotesize\color{obeGray}\code{nilai DW 1234h} \emph{tidak}
tersimpan sebagai \hex{12} lalu \hex{34}, melainkan \hex{34} lalu
\hex{12}; aturan yang sama berlaku untuk \dir{DD}. Urutan ini baru
terlihat bila memori dibaca \emph{byte per byte}.\par}
\Sumber{Bab IX, Urutan Little-Endian pada DW dan DD}
\end{frame}

% ---------------------------------------------------------
\begin{frame}{Rangkuman Bagian 1}
\small
\begin{itemize}
    \item Kode sumber memuat dua jenis baris: instruksi mesin yang
          menghasilkan opcode dan dieksekusi saat \eng{run time},
          serta direktif assembler (\eng{pseudo-op}) yang hanya
          berlaku saat \eng{assemble time} dan tidak menghasilkan
          kode mesin.

    \vspace{0.25em}
    \item \dir{DB} menyediakan 1 byte, \dir{DW} 2 byte, \dir{DD}
          4 byte, dan \dir{DT} 10 byte. Operator \dir{DUP} mengulang
          satu pola, sedangkan penanda \code{?} hanya menyediakan
          ruang tanpa nilai awal.

    \vspace{0.25em}
    \item Literal string tidak otomatis diakhiri; penanda \code{'\$'}
          ditambahkan sendiri untuk layanan DOS fungsi \hex{09}.
          Simbol \code{\$} sekaligus dapat dipakai menghitung
          \eng{offset}.

    \vspace{0.25em}
    \item Bilangan multi-byte disimpan \eng{little-endian}: byte
          berpangkat rendah berada di alamat lebih kecil, sehingga
          \code{DW 1234h} tersimpan \hex{34} lalu \hex{12}.
\end{itemize}
\end{frame}

% =========================================================
% BAGIAN 2
% =========================================================
\PemisahBagian{Direktif Simbol, Lokasi, dan Struktur Program}

% ---------------------------------------------------------
\begin{frame}{Konstanta Simbolik: \dir{EQU} dan \code{=}}
\footnotesize
\dir{EQU} memberi nama pada sebuah nilai. Setelah
\code{KELUAR EQU 4Ch} ditulis, setiap kemunculan \code{KELUAR}
digantikan \hex{4C} saat perakitan; nama semacam ini disebut
\term{konstanta simbolik}. Direktif \code{=} berbentuk serupa, yaitu
\code{URUTAN = 10}, tetapi sifatnya berbeda.

\vspace{0.35em}
\begin{center}
\begin{tabular}{@{}p{3.0cm}p{5.0cm}p{5.0cm}@{}}
    \toprule
    \textbf{Aspek} & \dir{EQU} & \code{=} \\
    \midrule
    Sifat simbol & Konstanta tetap sepanjang berkas &
    Nilai yang boleh berubah \\
    \addlinespace
    Dapat didefinisikan ulang & Tidak; assembler melaporkan
    \file{Symbol already defined} & Ya; definisi terakhir yang
    berlaku \\
    \addlinespace
    Jenis nilai yang boleh & Bilangan, alamat, teks, atau ekspresi
    yang dapat dihitung assembler & Wajib berupa bilangan \\
    \addlinespace
    Pemakaian khas & Nama nomor fungsi DOS, ukuran tabel &
    Penghitung di dalam makro atau blok yang diulang \\
    \bottomrule
\end{tabular}
\end{center}
\Sumber{Bab IX, Konstanta Simbolik}
\end{frame}

% ---------------------------------------------------------
\begin{frame}[fragile]{Definisi Ulang \dir{EQU} Dibandingkan \code{=}}
\footnotesize
Definisi ulang \dir{EQU} menimbulkan kesalahan perakitan:
\vspace{0.15em}
\begin{kodeasm}[numbers=none]
LEBAR   EQU 80          ; definisi pertama
LEBAR   EQU 40          ; **Error** Symbol already defined
\end{kodeasm}
{\color{obeGray}Pesan yang dikeluarkan TASM:\par}
\begin{keluaran}
**Error** equ-salah.asm(2) Symbol already defined: LEBAR
**Error** equ-salah.asm(2) Cannot define symbol
2 errors
\end{keluaran}

\vspace{0.1em}
Dengan \code{=}, definisi ulang menjadi sah dan nilai terakhir yang
berlaku:
\vspace{0.15em}
\begin{kodeasm}[numbers=none]
URUTAN  = 10            ; nilai pertama
URUTAN  = 20            ; sah: nilai sekarang menjadi 20
\end{kodeasm}
\Sumber{Bab IX, Definisi Ulang EQU}
\end{frame}

% ---------------------------------------------------------
\begin{frame}{Direktif Lokasi dan Struktur Program}
\footnotesize
\begin{description}
    \item[\dir{ORG}] Menetapkan penghitung lokasi
          (\eng{location counter}). Pada program \file{.COM},
          \code{ORG 100h} wajib ditulis karena DOS memuat program
          pada \eng{offset} \hex{100}.

    \vspace{0.3em}
    \item[\dir{SEGMENT} dan \dir{ENDS}] Menandai awal dan akhir
          sebuah segmen bernama. Setiap \dir{SEGMENT} harus ditutup
          \dir{ENDS} dengan nama segmen yang sama.

    \vspace{0.3em}
    \item[\dir{ASSUME}] Memberi tahu assembler register segmen mana
          yang dianggap menunjuk ke segmen tertentu, sehingga segmen
          \eng{default} sebuah label dapat ditentukan. Direktif ini
          tidak mengubah isi register saat program berjalan.

    \vspace{0.3em}
    \item[\dir{END}] Menutup berkas sumber; baris sesudahnya
          diabaikan. Dengan label, \code{END START} sekaligus
          mencatat titik masuk (\eng{entry point}). \dir{END} tanpa
          label lolos dari assembler, tetapi ditolak linker.

    \vspace{0.3em}
    \item[\dir{PROC} dan \dir{ENDP}] Menandai batas sebuah prosedur
          sehingga assembler dapat menentukan jenis pemanggilan.
          Keduanya tidak menghasilkan instruksi.
\end{description}
\Sumber{Bab IX, Direktif Lokasi dan Struktur}
\end{frame}

% ---------------------------------------------------------
\begin{frame}{Model Memori \dir{.MODEL}}
\footnotesize
\begin{center}
\begin{tabular}{@{}p{1.5cm}p{2.5cm}p{2.5cm}p{6.0cm}@{}}
    \toprule
    \textbf{Model} & \textbf{Batas kode} & \textbf{Batas data} &
    \textbf{Jumlah segmen} \\
    \midrule
    \file{TINY} & 64 KB & 64 KB &
    Satu segmen untuk kode, data, dan \eng{stack}; dipakai untuk
    \file{.COM} \\
    \addlinespace
    \file{SMALL} & 64 KB & 64 KB &
    Satu segmen kode dan satu segmen data; pilihan lazim untuk
    pembelajaran \\
    \addlinespace
    \file{MEDIUM} & Tidak dibatasi & 64 KB &
    Satu segmen data, banyak segmen kode \\
    \addlinespace
    \file{COMPACT} & 64 KB & Tidak dibatasi &
    Satu segmen kode, banyak segmen data \\
    \addlinespace
    \file{LARGE} & Tidak dibatasi & Tidak dibatasi &
    Banyak segmen; satu \eng{array} maksimum 64 KB \\
    \addlinespace
    \file{HUGE} & Tidak dibatasi & Tidak dibatasi &
    Seperti \file{LARGE}, tetapi \eng{array} boleh melewati 64 KB \\
    \bottomrule
\end{tabular}
\end{center}

\vspace{0.2em}
{\color{obeGray}\dir{.MODEL} juga menetapkan simbol bawaan seperti
\code{@DATA}. Semakin besar model, semakin banyak segmen yang harus
diurus program.\par}
\Sumber{Bab IX, Model Memori dan Struktur Prosedur}
\end{frame}

% ---------------------------------------------------------
\begin{frame}{Rangkuman Bagian 2}
\small
\begin{itemize}
    \item \dir{EQU} menghasilkan konstanta tetap yang tidak dapat
          didefinisikan ulang, sedangkan \code{=} menghasilkan nilai
          yang boleh diganti berkali-kali --- yang berlaku adalah
          definisi terakhir.

    \vspace{0.25em}
    \item Definisi ulang \dir{EQU} dilaporkan sebagai
          \file{Symbol already defined}; gunakan \code{=} bila
          memang membutuhkan nilai yang berubah.

    \vspace{0.25em}
    \item \dir{ORG} menetapkan penghitung lokasi (\code{ORG 100h}
          wajib pada \file{.COM}); \dir{SEGMENT}/\dir{ENDS} menandai
          batas segmen; \dir{ASSUME} menetapkan segmen
          \eng{default} bagi assembler; \dir{END} menutup berkas dan
          dapat menetapkan titik masuk.

    \vspace{0.25em}
    \item \dir{.MODEL} memilih model memori: \file{TINY} satu segmen
          untuk \file{.COM}, \file{SMALL} satu segmen kode dan satu
          segmen data. Model yang lebih besar memakai banyak segmen.

    \vspace{0.25em}
    \item \dir{PROC} dan \dir{ENDP} hanya menandai batas prosedur dan
          menentukan jenis pemanggilan, tanpa menghasilkan
          instruksi.
\end{itemize}
\end{frame}

% =========================================================
% BAGIAN 3
% =========================================================
\PemisahBagian{Peran Assembler dan Linker dalam Alur Translasi}

% ---------------------------------------------------------
\begin{frame}{Alur Translasi: \file{.asm} $\rightarrow$
\file{.obj} $\rightarrow$ \file{.exe}}
\small
\begin{center}
\begin{tikzpicture}[
    berkas/.style={draw=obeBlue, line width=0.8pt, fill=obeBlueLight,
        rounded corners=3pt, minimum width=1.9cm, minimum height=0.8cm,
        align=center, font=\scriptsize\bfseries, inner sep=2pt},
    alat/.style={berkas, draw=obeTeal, fill=obeTealLight},
    hasil/.style={berkas, draw=obeAmber, fill=obeAmberLight},
    panah/.style={-{Stealth[length=2mm]}, line width=0.9pt, obeTeal},
    panel/.style={draw=obeGray, line width=0.6pt, dashed,
        rounded corners=2pt, fill=obeGrayLight, align=left,
        font=\ttfamily\tiny, inner sep=3pt, text=black,
        text width=2.25cm},
]
    \node[berkas] (src) {\file{program.asm}};
    \node[alat, right=0.4cm of src] (asm) {ASSEMBLER};
    \node[hasil, right=0.4cm of asm] (obj) {\file{program.obj}};
    \node[alat, right=0.4cm of obj] (lnk) {LINKER};
    \node[hasil, right=0.4cm of lnk] (exe) {\file{program.exe}};

    \draw[panah] (src) -- (asm);
    \draw[panah] (asm) -- (obj);
    \draw[panah] (obj) -- (lnk);
    \draw[panah] (lnk) -- (exe);

    \node[panel, below=0.40cm of asm, anchor=north] (simbol) {
        \textbf{TABEL SIMBOL}\\
        START\ \ 0000H\\
        panjang\ 0012H
    };
    \node[panel, below=0.40cm of lnk, anchor=north] (relok) {
        \textbf{CATATAN RELOKASI}\\
        +0009\ OFF\ sapaan\\
        +0036\ SEG\ \_STACK
    };
\end{tikzpicture}
\end{center}

\vspace{0.3em}
{\footnotesize\color{obeGray}Assembler membangun tabel simbol dan
menghasilkan kode mesin; linker menyelesaikan \eng{relocation} dan
menentukan alamat akhir.\par}
\Sumber{Bab IX, Alur translasi .asm .obj .exe}
\end{frame}

% ---------------------------------------------------------
\begin{frame}{Tugas Assembler: Dua Lintasan}
\small Assembler bekerja dalam dua lintasan (\eng{two-pass}):
\vspace{0.3em}
\begin{enumerate}
    \item \textbf{Lintasan pertama.} Assembler membaca seluruh baris,
          memeriksa sintaks, menghitung besar setiap direktif data,
          dan mencatat setiap label beserta \eng{offset}-nya ke dalam
          \term{tabel simbol}. Pada lintasan ini belum ada kode mesin
          yang ditulis ke berkas.

    \vspace{0.45em}
    \item \textbf{Lintasan kedua.} Assembler menerjemahkan setiap
          instruksi menjadi opcode, menggabungkannya dengan
          \eng{offset} dari tabel simbol, lalu menuliskan hasilnya ke
          berkas objek. Karena itu, label yang terletak
          \emph{setelah} instruksi pemakainya tetap dapat dipakai.
\end{enumerate}

\vspace{0.45em}
\begin{block}{Berkas \file{.lst}}
\small Berkas \file{.lst} memperlihatkan hasil kedua lintasan
sekaligus: kolom kode mesin, tabel simbol, dan \eng{offset} yang sudah
dihitung.
\end{block}
\Sumber{Bab IX, Tugas Assembler}
\end{frame}

% ---------------------------------------------------------
\begin{frame}[fragile]{Membaca Tabel Simbol pada Berkas \file{.lst}}
\footnotesize
\begin{keluaran}
  4                             .DATA
  5  0000  48 61 6C 6F 2C 20 61 73 73 65  sapaan  DB 'Halo, ...'
  7  0012                        panjang EQU $ - sapaan
  8  0012  41 42 43 44           kata    DW 4241h, 4443h
  9  0016  78 56 34 12           nilai   DD 12345678h
 10  001A  2A 2A 2A 2A           bintang DB 4 DUP('*')
 11  001E  ??                    buffer  DB 8 DUP(?)
 14  0000  B8 ----            START:  MOV AX, @DATA
 17  0007  8D 16 0000 R               LEA DX, sapaan
\end{keluaran}

\vspace{0.2em}
{\color{obeGray}
\begin{itemize}
    \item Direktif data terlihat memakan ruang memori, sedangkan baris
          \dir{.MODEL}, \dir{.STACK}, dan \dir{.CODE} tidak.
    \item \code{panjang} tercatat \hex{12}, yaitu \dec{18}, hasil
          perhitungan \code{\$ - sapaan} oleh assembler.
    \item Tanda \code{R} berarti baris itu memuat \term{relokasi}:
          alamatnya masih menunggu linker.
\end{itemize}\par}
\Sumber{Bab IX, Cuplikan direktif-data.lst}
\end{frame}

% ---------------------------------------------------------
\begin{frame}{Tugas Linker: Resolusi Simbol dan \eng{Relocation}}
\small Berkas \file{.obj} belum dapat dijalankan karena masih memuat
\term{simbol eksternal}, yaitu nama yang belum punya alamat pasti.
Linker melakukan tiga pekerjaan pokok:
\vspace{0.3em}
\begin{enumerate}
    \item \textbf{Resolusi simbol eksternal.} Setiap rujukan diberi
          alamat sesungguhnya. Contoh yang paling akrab adalah
          \code{@DATA}, yang baru memperoleh nomor segmen setelah
          linker menempatkan segmen data.

    \vspace{0.4em}
    \item \textbf{\eng{Relocation}.} Linker mencatat bagian kode mana
          saja yang harus disesuaikan bila program dimuat pada alamat
          dasar yang berbeda.

    \vspace{0.4em}
    \item \textbf{Penentuan alamat akhir.} Linker menyusun seluruh
          segmen berurutan, menetapkan alamatnya, dan menggabungkan
          berkas objek dengan pustaka (\eng{library}) bila
          diperlukan.
\end{enumerate}

\vspace{0.35em}
{\footnotesize\color{obeGray}Ringkasan pekerjaan itu terekam pada
berkas \file{.map} yang dibuat dengan opsi \file{/m}.\par}
\Sumber{Bab IX, Tugas Linker}
\end{frame}

% ---------------------------------------------------------
\begin{frame}[fragile]{Berkas \file{.map}: Alamat Akhir Setiap Segmen}
\footnotesize
\begin{keluaran}
Turbo Link  Version 5.0.  Copyright (c) 1992 Borland International

 Start  Stop   Length Name                   Class
 00000H 00035H 00036H _TEXT                  CODE
 00036H 0005BH 00026H _DATA                  DATA
 0005CH 0005BH 00000H _BSS                   BSS
 0005CH 00015BH 00100H _STACK                 STACK

Program entry point at 0000:0000
\end{keluaran}

\vspace{0.2em}
{\color{obeGray}Segmen \file{\_DATA} berukuran \hex{26} byte, tepat
sama dengan total direktif data pada berkas sumber. Titik masuk
tercatat \file{0000:0000}, yaitu label \file{START}.\par}
\Sumber{Bab IX, Cuplikan direktif-data.map}
\end{frame}

% ---------------------------------------------------------
\begin{frame}[fragile]{Alur Lengkap dari Baris Perintah DOS}
\footnotesize
\begin{keluaran}
C:\ASM>tasm /l direktif-data.asm
Turbo Assembler  Version 4.1
Assembling file:   direktif-data.asm
Error messages:    None
Passes:            2

C:\ASM>tlink /m direktif-data.obj
Turbo Link  Version 5.0.

C:\ASM>direktif-data
Halo, assembly!
ABCD****
\end{keluaran}

\vspace{0.2em}
{\color{obeGray}Assembler menghasilkan \file{.obj} dan --- karena opsi
\file{/l} --- \file{.lst}; linker menghasilkan \file{.exe} dan ---
karena opsi \file{/m} --- \file{.map}. Keluaran \code{ABCD****}
terjadi karena data dibaca \emph{byte per byte} secara
\eng{little-endian}.\par}
\Sumber{Bab IX, Alur lengkap dari DOS}
\end{frame}

% ---------------------------------------------------------
\begin{frame}{Empat Kesalahan Perakitan yang Paling Sering}
\footnotesize
\begin{enumerate}
    \item \textbf{Simbol ganda.} Satu nama dipakai untuk dua label
          atau dua \dir{EQU}. Pesan \file{Symbol already defined}
          hanya menyebut nama yang bentrok, sehingga letaknya harus
          dicari pada berkas \file{.lst}.

    \vspace{0.3em}
    \item \textbf{\dir{EQU} didefinisikan ulang.} Pesan
          \file{Symbol already defined} pada baris \dir{EQU} kedua.
          Gunakan \code{=} bila memang membutuhkan nilai yang
          berubah.

    \vspace{0.3em}
    \item \textbf{\dir{END} tanpa label.} Perakitan selesai tanpa
          kesalahan, tetapi linker melaporkan tidak menemukan titik
          masuk program. Tuliskan \code{END START} sesuai nama label
          awal.

    \vspace{0.3em}
    \item \textbf{Segmen tidak ditutup \dir{ENDS}.} Muncul pesan
          \file{Block nesting error} atau \file{Missing ENDS}, dan
          semua baris sesudahnya ikut salah tafsir.
\end{enumerate}
\Sumber{Bab IX, Kesalahan perakitan yang umum}
\end{frame}

% ---------------------------------------------------------
\begin{frame}{Rangkuman Bagian 3}
\small
\begin{itemize}
    \item Assembler bekerja dua lintasan: lintasan pertama menghitung
          ukuran direktif dan membangun tabel simbol, lintasan kedua
          menghasilkan opcode dengan memakai \eng{offset} dari tabel
          itu.

    \vspace{0.25em}
    \item Berkas \file{.lst} memperlihatkan hasil kedua lintasan,
          termasuk tanda \code{R} yang berarti baris tersebut masih
          menunggu \eng{relocation} oleh linker.

    \vspace{0.25em}
    \item Linker menyelesaikan simbol eksternal (misalnya
          \code{@DATA}), melakukan \eng{relocation}, dan menetapkan
          alamat akhir; ringkasannya terekam pada berkas \file{.map}.

    \vspace{0.25em}
    \item Alur lengkap: \file{tasm /l program.asm} $\rightarrow$
          \file{tlink /m program.obj} $\rightarrow$ menjalankan
          \file{program.exe}.

    \vspace{0.25em}
    \item Empat kesalahan tersering: simbol ganda, \dir{EQU}
          didefinisikan ulang, \dir{END} tanpa label, dan segmen
          tanpa \dir{ENDS}.
\end{itemize}
\end{frame}

% =========================================================
% CONTOH PROGRAM
% =========================================================
\begin{frame}[fragile]{Contoh 1: Mendemonstrasikan Semua Direktif
Data}
\footnotesize
\begin{columns}[t]
\begin{column}{0.54\textwidth}
\begin{kodeasm}[numbers=none]
.MODEL SMALL
.STACK 100h
.DATA
    sapaan  DB 'Halo, assembly!', 0Dh, 0Ah, '$'
    panjang EQU $ - sapaan       ; 18 byte
    kata    DW 4241h, 4443h
    nilai   DD 12345678h
    bintang DB 4 DUP('*')
    buffer  DB 8 DUP(?)
\end{kodeasm}
\end{column}
\begin{column}{0.44\textwidth}
\begin{kodeasm}[numbers=none]
.CODE
START:
    MOV AX, @DATA
    MOV DS, AX
    MOV AH, 09h
    LEA DX, sapaan
    INT 21h
    MOV DL, BYTE PTR [kata]
    INT 21h
    MOV AH, 4Ch
    INT 21h
END START
\end{kodeasm}
\end{column}
\end{columns}

\vspace{0.15em}
{\color{obeTeal}\bfseries Keluaran}\quad
\begin{keluaran}
Halo, assembly!
ABCD****
\end{keluaran}
{\color{obeGray}Data dibaca \emph{byte per byte} dengan \code{BYTE PTR};
baris kedua keluaran membuktikan urutan \eng{little-endian}.\par}
\Sumber{Bab IX, Contoh 1 --- \file{code/asm/direktif-data.asm}}
\end{frame}

% ---------------------------------------------------------
\begin{frame}[fragile]{Contoh 2: Konstanta \dir{EQU}}
\footnotesize
\begin{columns}[t]
\begin{column}{0.50\textwidth}
\begin{kodeasm}[numbers=none]
LAYAR      EQU 02h    ; cetak karakter
TULIS      EQU 09h    ; cetak string
KELUAR     EQU 4Ch    ; akhiri program
BARIS_BARU EQU 0Dh    ; carriage return
AWAL_BARIS EQU 0Ah    ; line feed

.DATA
  judul  DB 'Konstanta EQU', 0Dh, 0Ah, '$'
  isi    DB 'X = ', '$'
\end{kodeasm}

\vspace{0.2em}
{\color{obeTeal}\bfseries Keluaran}\quad
\begin{keluaran}
Konstanta EQU
X = 7
\end{keluaran}
\end{column}
\begin{column}{0.47\textwidth}
\begin{kodeasm}[numbers=none]
.CODE
START:
    MOV AX, @DATA
    MOV DS, AX
    MOV AH, TULIS        ; cetak judul
    LEA DX, judul
    INT 21h
    MOV AH, TULIS        ; cetak 'X = '
    LEA DX, isi
    INT 21h
    MOV AH, LAYAR        ; cetak '7'
    MOV DL, '7'
    INT 21h
    MOV AH, KELUAR
    INT 21h
END START
\end{kodeasm}
\end{column}
\end{columns}
\Sumber{Bab IX, Contoh 2 --- \file{code/asm/equ-konstanta.asm}}
\end{frame}

% =========================================================
% PRAKTIKUM
% =========================================================
\begin{frame}{Praktikum: Merakit dengan TASM dan TLINK}
\begin{block}{Judul praktikum}
\footnotesize Merakit dengan TASM dan TLINK serta memeriksa berkas
\file{.lst} dan \file{.map} \quad\textbullet\quad
\textbf{Sub-CPMK:} 2.3 \quad\textbullet\quad
\textbf{Waktu:} 100 menit \quad\textbullet\quad
\textbf{Bentuk:} individu
\end{block}

\vspace{0.2em}
\begin{columns}[t]
\begin{column}{0.46\textwidth}
{\color{obeBlue}\bfseries Alat dan bahan}
\begin{itemize}\footnotesize
    \item DOSBox dengan direktori kerja ter-\eng{mount}
    \item TASM dan TLINK versi DOS
    \item \file{direktif-data.asm} dan \file{equ-konstanta.asm}
    \item Lembar kerja pengamatan
\end{itemize}
\end{column}
\begin{column}{0.50\textwidth}
{\color{obeBlue}\bfseries Langkah kerja}
\begin{enumerate}\footnotesize
    \item \file{tasm /l direktif-data.asm}; catat lintasan dan
          kesalahan
    \item \file{tlink /m direktif-data.obj}; jalankan \file{.exe}
    \item Periksa \file{.lst}: baris direktif dan baris bertanda
          \code{R}
    \item Periksa \file{.map}: hitung besar \file{\_DATA} dan
          \file{\_TEXT}
    \item Sengaja buat kesalahan \dir{EQU} dan \dir{END} tanpa
          label
\end{enumerate}
\end{column}
\end{columns}

\vspace{0.2em}
{\footnotesize\textbf{Dikumpulkan:} berkas \file{.asm} beserta
\file{.obj}, \file{.lst}, \file{.map}, dan \file{.exe}, serta catatan
kesalahan yang disengaja dan cara memperbaikinya.\par}
\end{frame}

% =========================================================
% LATIHAN
% =========================================================
\begin{frame}{Latihan}
\small
\begin{enumerate}
    \item Jelaskan perbedaan instruksi mesin dan direktif assembler,
          serta mengapa direktif disebut \eng{pseudo-op}. Berikan dua
          contoh direktif yang mengisi memori dan dua contoh yang
          tidak mempengaruhi isi berkas \file{.exe}.

    \vspace{0.25em}
    \item Sebutkan direktif data yang tepat, lalu tuliskan
          deklarasinya: (a) satu nilai \hex{0A} untuk pindah baris;
          (b) tabel delapan nilai 16-bit \hex{0000}; (c) penampung 32
          byte yang diisi saat program berjalan; (d) string
          \code{'SEMARANG'} yang siap dicetak fungsi \hex{09}.

    \vspace{0.25em}
    \item Diberikan blok data \code{a DB 'TEST'}, \code{b DW 2 DUP(0)},
          \code{c DD 1, 2}, \code{d DB 3 DUP(?)}, dan \code{e DT 0}.
          Hitung total byte yang disediakan dan tuliskan \eng{offset}
          setiap label.

    \vspace{0.25em}
    \item Tuliskan urutan byte di memori untuk \code{DW 0102h} dan
          \code{DD 03040506h}, lalu jelaskan mengapa urutannya
          demikian.

    \vspace{0.25em}
    \item Sebuah berkas sumber berakhir dengan \code{END} tanpa nama
          label. Apakah TASM melaporkan kesalahan? Apakah TLINK
          melaporkan kesalahan? Bagaimana cara memperbaikinya?
\end{enumerate}
\end{frame}

% =========================================================
% RANGKUMAN
% =========================================================
\begin{frame}{Rangkuman}
\small
\begin{itemize}
    \item Kode sumber assembly terdiri atas dua jenis baris:
          instruksi mesin yang menghasilkan opcode dan dieksekusi saat
          program berjalan, serta direktif assembler yang hanya
          berlaku saat perakitan dan tidak menghasilkan kode mesin.

    \vspace{0.25em}
    \item \dir{DB} menyediakan satu byte, \dir{DW} dua byte,
          \dir{DD} empat byte, dan \dir{DT} sepuluh byte. \dir{DUP}
          mengulang satu pola, sedangkan \code{?} hanya menyediakan
          ruang tanpa nilai awal.

    \vspace{0.25em}
    \item Bilangan multi-byte disimpan \eng{little-endian}, sehingga
          \code{DW 1234h} tersimpan sebagai byte \hex{34} lalu
          \hex{12}.

    \vspace{0.25em}
    \item \dir{EQU} membuat konstanta tetap yang tidak dapat
          didefinisikan ulang, sedangkan \code{=} membuat nilai yang
          boleh diganti berkali-kali.

    \vspace{0.25em}
    \item \dir{ORG}, \dir{SEGMENT}/\dir{ENDS}, \dir{ASSUME}, dan
          \dir{END} mengatur letak dan struktur program; \dir{END}
          yang diberi label sekaligus menetapkan titik masuk.

    \vspace{0.25em}
    \item Assembler membangun tabel simbol dan menghitung
          \eng{offset}; linker menyelesaikan simbol eksternal,
          melakukan \eng{relocation}, dan merekam alamat akhir pada
          berkas \file{.map}.
\end{itemize}
\end{frame}

% ---------------------------------------------------------
\begin{frame}{Cek Pemahaman}
\small
\begin{enumerate}
    \item Pernyataan yang benar tentang direktif assembler adalah
          \dots
          \begin{enumerate}\footnotesize
              \item direktif menghasilkan opcode seperti instruksi
              \item direktif hanya berlaku saat perakitan dan tidak
                    menghasilkan kode mesin
              \item direktif dieksekusi prosesor saat program berjalan
              \item direktif tidak berpengaruh pada berkas \file{.exe}
          \end{enumerate}

    \vspace{0.25em}
    \item Deklarasi \code{nilai DW 1234h} menempatkan byte di memori
          dengan urutan \dots
          \begin{enumerate}\footnotesize
              \item \hex{12} lalu \hex{34} \quad
              \item \hex{34} lalu \hex{12} \quad
              \item \hex{00} lalu \hex{12} \quad
              \item \hex{43} lalu \hex{21}
          \end{enumerate}

    \vspace{0.25em}
    \item Tugas yang dikerjakan linker, bukan assembler, adalah
          \dots
          \begin{enumerate}\footnotesize
              \item menghitung \eng{offset} label dalam satu segmen
              \item menerjemahkan mnemonic menjadi opcode
              \item menyelesaikan simbol eksternal dan melakukan
                    \eng{relocation}
              \item membaca direktif \dir{EQU}
          \end{enumerate}
\end{enumerate}
\end{frame}

% =========================================================
% PENUTUP
% =========================================================
\begin{frame}{Penutup}
\small
\begin{block}{Sebelum pertemuan berikutnya}
\begin{itemize}
    \item Selesaikan Tugas 3: berkas hasil perakitan, cuplikan
          \file{.lst} dan \file{.map}, serta catatan kesalahan
          perakitan yang disengaja.
    \item Tuliskan deklarasi \dir{DB}, \dir{DW}, dan \dir{DD} untuk
          tabel lima nilai, lalu hitung \eng{offset} setiap label.
    \item Baca \textbf{Bab X} --- Instruksi Transfer Data dan
          Aritmatika.
\end{itemize}
\end{block}

\vspace{0.5em}
\begin{center}
{\color{obeBlue}\sffamily\large\bfseries Terima kasih\par}
\vspace{0.2em}
{\color{obeGray}\sffamily\small \NamaMK\ (\KodeMK) \textbullet\
\SemesterTeks\par}
\end{center}
\end{frame}

\end{document}
